\section{The inverse-sumset route to upper bounds}
\label{sec:inverse-sumset-route}

For an \(m\)-sum-free set in \(\Zp\), with \(p\nmid m\), the sets \(2A\) and
\(mA\) are disjoint.  Since multiplication by \(m\) is a bijection,
\[
    |2A|+|A|\leq p.
\]
Together with Cauchy--Davenport, \(|2A|\geq2|A|-1\), this gives the baseline
upper bound \(|A|\leq(p+1)/3\).  A density close to \(1/3\) forces \(2A\) to be
small, so an inverse theorem for small sumsets can place \(A\) inside a short
arithmetic progression.  One can then use the disjointness of \(2A\) and
\(mA\) again, now with the extra progression structure, to improve the bound.

This is the route developed on the circle by Candela and de Roton
\cite{CandelaDeRoton2019}.  Their continuous small-doubling theorem, using the
then available Serra--Z\'emor constant \(10^{-4}\), yields
\[
    d_m(\T)\leq \frac{1}{3.0001}
\]
for all \(m\geq3\).  Candela, Gonz\'alez-S\'anchez, and Grynkiewicz subsequently
proved a finite-group inverse result optimized for the density range of this
application and obtained the stronger uniform estimate
\[
    d_m(\T)
    =\lim_{p\to\infty}d_m(\Zp)
    \leq \frac1{3.1955}
\]
\cite[Theorem~1.5]{CandelaGonzalezSanchezGrynkiewicz2021}.  These works build on
the prime-cyclic \(3k-4\) program, including the theorem of Serra and Z\'emor
\cite{SerraZemor2009}.

\subsection{What the 2024 theorem contributes}

Grynkiewicz proves new high-density forms of the \(3k-4\) theorem for general
sumsets \(A+B\) in \(\Zp\) \cite{Grynkiewicz2024}.  Two features are relevant
here: tangible small-doubling constants in a high-density regime, and a result
under the conjecturally ideal density restriction.  The hypotheses are not
identical to those used to prove the bound \(1/3.1955\).  Therefore the paper is
an input to a new calculation, not by itself a new \(m\)-sum-free density
bound.

A concrete next task is to repeat the optimization in
\cite[Section~3]{CandelaGonzalezSanchezGrynkiewicz2021} with the strongest
applicable 2024 inverse statement.  Write
\[
    |A|=\eta p,
    \qquad
    |2A|=2|A|+r,
    \qquad
    3|A|+r\leq p.
\]
For every branch of the inverse theorem one must:
\begin{enumerate}
    \item translate its hypotheses into inequalities in \(\eta\) and \(r/p\);
    \item determine the density range in which those hypotheses are automatic;
    \item insert the resulting progression lengths into the intersection
          argument for \(2A\) and \(mA\); and
    \item optimize the final one-variable inequality in \(\eta\).
\end{enumerate}
This calculation should be written independently of any numerical interval
model: it concerns arbitrary subsets of \(\Zp\).

\subsection{The gap as \texorpdfstring{$m\to\infty$}{m tends to infinity}}

The construction in Section~\ref{sec:equally-spaced-construction} gives
\[
    d_m(\T)\geq
    \frac{\lfloor(m+2)^2/8\rfloor}{m(m+2)}
    =\frac18+\frac{1}{4m}+O\!\left(\frac1{m^2}\right),
\]
up to the bounded rounding term in the numerator.  The uniform upper bound
\(1/3.1955\) does not approach this scale.  There are therefore two distinct
upper-bound projects:
\begin{itemize}
    \item improve the best uniform constant by importing stronger inverse
          sumset results; and
    \item obtain an \(m\)-dependent argument that uses the large multiplier,
          with the eventual target \(d_m(\T)\leq1/8+o(1)\).
\end{itemize}

The second goal needs information that the coarse inequality
\(\mu(2A)+\mu(mA)\leq1\) discards.  Promising quantities to retain are the
number and arrangement of components of \(m^{-1}(2A)\), the distribution of
\(A\) among arcs of length \(1/m\), and the interaction between a progression
containing \(A\) and its dilation by \(m\).  At present these are research
directions, not established reductions.
